% ============================================================
%  chapters/ch12-keyboards-gesture-sensors.tex
% ============================================================

\chapter{Keyboards as Gesture Sensors}

\chapterepigraph{%
  Every key you press leaves a signature\\
  in kinematic space that the text does not record.\\
  The keyboard hears more than it reports.
}{Flyxion, \textit{Semantic Infrastructure}}

\sidenote{App.\,G for the\\full keyboard\\geometry\\formalism}

The keyboard is the primary input device of the information age. It has
been analyzed as a symbol source (by information theorists), as an
ergonomic tool (by human factors engineers), and as a cultural artifact
(by historians of technology). This chapter analyzes it as something
different: a \emph{gesture sensor} — a device that samples the continuous
kinematic behavior of a person's hands and reports a discrete projection
of that behavior.

\section{Typing as Coordinated Motor Behavior}

A trained typist's hands move continuously, with overlapping finger
activations, rhythmic patterns, and smooth inter-key transitions.
The keyboard, however, reports only discrete events: key-down and key-up,
time-stamped to millisecond precision.

\begin{definition}{Typing Stream}{typing-stream}
  A \emph{typing stream} is the sequence of keyboard events:
  \[
    T = \bigl((k_1, t_1^\text{down}, t_1^\text{up}),
               (k_2, t_2^\text{down}, t_2^\text{up}), \ldots\bigr)
  \]
  where $k_i \in K$ is a key and $t_i^\text{down} < t_i^\text{up}$
  are the key-press and key-release times.
\end{definition}

The typing stream is the \emph{projection} of the full kinematic
hand trajectory onto the discrete key-event space. The projection
discards:

\begin{itemize}
  \item Finger position between key presses (the flight phase),
  \item Finger velocity and acceleration (the dynamics),
  \item Inter-finger coordination (simultaneous movements),
  \item Hand position and wrist posture,
  \item Pressure beyond the binary press/release threshold.
\end{itemize}

\section{The Hidden Signal Beneath Text}

The typing stream contains considerably more information than the
text it produces, and the text contains considerably less than the
typing stream.

\begin{definition}{Keystroke Dynamics}{keystroke-dynamics}
  The \emph{keystroke dynamics} of a typing session are the timing
  features of the typing stream:
  \begin{itemize}
    \item \textbf{Dwell time:} $D_i = t_i^\text{up} - t_i^\text{down}$
          (how long each key is held).
    \item \textbf{Flight time:} $F_i = t_{i+1}^\text{down} - t_i^\text{up}$
          (time between key releases and next key presses).
    \item \textbf{Digraph latency:} $L_{ij} = t_j^\text{down} - t_i^\text{down}$
          (time between successive key presses).
    \item \textbf{Overlap duration:} $O_{ij} = \min(t_i^\text{up}, t_j^\text{up})
          - t_j^\text{down}$ (duration of simultaneous key presses).
  \end{itemize}
\end{definition}

These timing features are individual-specific, consistent across sessions,
and robust to training. They constitute a \emph{biometric}: keystroke
dynamics can identify individuals with accuracy comparable to fingerprinting.
The text they produce contains none of this information.

\begin{conceptaside}{Text as Information Bottleneck}
  The conventional view of typing: person $\to$ text $\to$ reader.

  The gesture view: person $\to$ kinematic stream $\to$ typing stream
  $\to$ text $\to$ reader.

  Each arrow is a lossy projection. The text at the end contains
  only the symbolic content. The kinematic stream contains the
  motor signature of the person. The typing stream is between them:
  richer than text (it includes timing), poorer than kinematics
  (it discards position).

  Current AI language models are trained on text — the end of the
  chain. They have no access to the kinematic or timing information.
  This is why they cannot detect fatigue, emotion, or identity from
  text alone: that information was discarded two projections upstream.
\end{conceptaside}

\section{Overlap, Pivots, and Release Structures}

The most information-rich features of the typing stream are those
that encode the coordination structure of the two hands.

\begin{definition}{Key Overlap}{key-overlap}
  Keys $k_i$ and $k_j$ \emph{overlap} if their press intervals intersect:
  \[
    [t_i^\text{down}, t_i^\text{up}] \cap [t_j^\text{down}, t_j^\text{up}]
    \neq \emptyset.
  \]
  The \emph{overlap pattern} of a word is the set of overlapping key pairs.
\end{definition}

Overlap patterns are kinematically determined: which keys overlap depends
on which fingers are used and how the transitions between them are timed.
High-skilled typists show characteristic overlap patterns that low-skilled
typists do not — the patterns are signatures of motor expertise.

\begin{definition}{Typing Pivot}{typing-pivot}
  A \emph{typing pivot} occurs at position $i$ in a word when the
  typing changes hands:
  \[
    H(k_{i-1}) \neq H(k_i)
  \]
  where $H(k) \in \{L, R\}$ is the hand assignment.
  Pivots are characterized by a brief acceleration (flight time
  $F_i$ is typically shorter at pivots due to parallel hand preparation).
\end{definition}

\begin{definition}{Release Structure}{release-struct}
  The \emph{release structure} of a word is the ordering of key-up events:
  whether key $i$ is released before or after key $j$ is pressed.
  Release structures encode the chunking of words into motor programs.
\end{definition}

\section{Inter-Hand Coordination}

The deepest structure in typing kinematics is the coordination between
the two hands. Because the two hands are semi-independent motor systems,
their coordination creates rich temporal structure.

\begin{definition}{Hand Coordination Profile}{hand-coord}
  The \emph{hand coordination profile} of a typing session is the
  function $\rho : K \times K \to \mathbb{R}$ where $\rho(k_i, k_j)$
  is the average flight time from $k_i$ to $k_j$ across all occurrences.
  Same-hand transitions have $\rho(k_i,k_j) > 0$; cross-hand transitions
  tend to have smaller $\rho$ due to parallel preparation.
\end{definition}

The hand coordination profile is a complete characterization of the
timing structure of a person's typing. Two typists producing the same
text at the same overall speed can have radically different coordination
profiles — reflecting different motor programs for the same symbolic output.

\section{Keyboard Geometry Revisited}

The keyboard geometry formalism of Appendix~G defines keys as vertices
in a motor graph $\Gamma = (K, E)$. The typing stream is a walk on
this graph. The walk's properties — edge traversal times, node dwell
times, vertex degree sequences — encode the kinematic structure of
the typing.

\begin{theorem}{Keyboard Graph is a Gesture Sensor}{keyboard-sensor}
  The keyboard graph $\Gamma$ and the typing stream $T$ together
  determine the first-order kinematic structure of the typing gesture:
  specifically, the digraph latency profile $\{L_{ij}\}$ and the
  overlap pattern.
\end{theorem}

\begin{proof}
  The digraph latency $L_{ij} = t_j^\text{down} - t_i^\text{down}$ is
  directly available from $T$. The overlap pattern is determined by
  comparing press/release intervals. Both are fully determined by $T$.
  The keyboard graph $\Gamma$ provides the spatial embedding that
  interprets the timing: knowing which keys are adjacent, which hand
  they belong to, and their stretch costs allows the latencies to be
  decomposed into flight phase, preparation, and keystroke components.
\end{proof}

\begin{exercise}
  For the word ``semantic'', identify all key pairs that overlap
  (assume standard QWERTY layout and 10-finger touch typing).
  Where are the typing pivots? Describe the expected release structure.
\end{exercise}

\begin{exercise}
  Define the \emph{typing entropy} of a word $w$ as the entropy of
  the distribution of possible overlap patterns over skilled typists
  of $w$. Is typing entropy higher for frequently typed words
  or rare words? Explain intuitively why.
\end{exercise}

\begin{exercise}[title={Research}]
  The digraph latency profile $\{L_{ij}\}$ is a matrix of size
  $|K| \times |K|$. For standard QWERTY, $|K| \approx 47$.
  This matrix has $\approx 2200$ entries, most of which are rarely
  observed in practice. Propose a dimensionality-reduction technique
  (PCA, manifold learning, or another method) that extracts the
  dominant modes of variation in $\{L_{ij}\}$ across typists.
  What do you expect the dominant modes to represent?
\end{exercise}
